\section{Derivations}\label{app:derivations}

Let the AI-induced annual increment at date $\tau$ be
$x_{i,s,\tau}=\phi^{Y,nom}_{s,\tau}T_{i,s,\tau}$. Starting from the no-AI
counterfactual level $Y_i^0$, compounding gives
\begin{equation}
  Y^{AI}_{i,s,H}
  =Y_i^0\prod_{\tau=1}^{H}(1+x_{i,s,\tau}),
  \qquad
  g^{AI}_{i,s,H}
  =\frac{Y^{AI}_{i,s,H}-Y_i^0}{Y_i^0}
  =\prod_{\tau=1}^{H}
    \left(1+\phi^{Y,nom}_{s,\tau}T_{i,s,\tau}\right)-1.
  \label{eq:accumulated_gain_derivation}
\end{equation}
When the annual scenario and domestic translation are constant over the
horizon, this reduces to
$g^{AI}=(1+\phi^{Y,nom}T)^H-1$.

Suppress indices and write
$\widetilde M=MFC^{gross}-D^{ME}$. Substituting
$fs^{eff}=\widetilde M g^{AI}-F^{fix}$ into the buffered accounting
condition yields
\begin{align}
  fs^{eff}&\geq(1+\xi)c^{gross}
  \nonumber\\
  \Longleftrightarrow\quad
  \widetilde M g^{AI}
    &\geq(1+\xi)c^{gross}+F^{fix}
     \equiv\mathcal R(\xi).
  \label{eq:accounting_requirement_derivation}
\end{align}
The debt guardrail adds the inherited primary-balance gap to the same
requirement:
\begin{equation}
  \widetilde M g^{AI}
  \geq \mathcal R(\xi)+s^{PB}
  \equiv \mathcal R^{DC}(\xi).
  \label{eq:debt_requirement_derivation}
\end{equation}
Thus fixed deductions enter the requirement once, while marginal-equivalent
deductions enter the capture rate once. The corresponding unbuffered ratio is
$V^{gross}=(\widetilde M g^{AI}-F^{fix})/c^{gross}$; setting $V^{gross}=1$
is Equation~\eqref{eq:accounting_requirement_derivation} at $\xi=0$.

Let $\mathcal Q$ denote either $\mathcal R$ or $\mathcal R^{DC}$. At the
boundary $\widetilde M g^{AI}=\mathcal Q$, direct rearrangement gives
\begin{align}
  g^{AI,*}
    &=\frac{\mathcal Q}{\widetilde M},
  &
  MFC^{gross,*}
    &=\frac{\mathcal Q}{g^{AI}}+D^{ME}.
  \label{eq:gain_capture_derivation}
\end{align}
For a constant annual scenario and fixed domestic structure, substituting
$g^{AI}=(1+\phi^{Y,nom}T)^H-1$ and solving the same boundary for translation
gives
\begin{equation}
  T^{req,H}
  =\frac{\left(1+\mathcal Q/\widetilde M\right)^{1/H}-1}
          {\phi^{Y,nom}}.
  \label{eq:translation_requirement_derivation}
\end{equation}
These closed forms require $g^{AI}>0$, $\widetilde M>0$, and
$\phi^{Y,nom}>0$ for the corresponding inversion. A non-positive denominator
does not produce a finite positive crossing. Replacing $\mathcal R$ with
$\mathcal R^{DC}$ gives the debt-consistent inversion without changing the
algebra. If channel weights or endpoint costs are recomputed as the inverted
quantity changes, the boundary is instead solved numerically.
